\section{Finite Clebsch covers and rigidity}\label{sec:clebsch-finite}

The preceding cover bounds are asymptotic and permit varying coefficients.
Here we fix the ordinary adjacency check and obtain an exact finite bound
without a separating-chain hypothesis. These are different restrictions;
neither set of hypotheses is supplied by the other.

Let $G$ be the Clebsch graph on $\F_2^4$, identified with the integers
$0,\ldots,15$, where $u\sim v$ when $u\mathbin\oplus v$ belongs to
$\{1,2,4,8,15\}$. Write $M=I+A_G$.
Equivalently, label the five generators by $i\in[5]$ and their ten
pairwise sums by $ij$. The edges are $0$--$i$, $i$--$ij$, and disjoint
pairs. Over any characteristic-two field,
\begin{equation}\label{eq:clebsch-base}
 \rank M=6,\qquad \dim\ker M=10,\qquad M^2=0.
\end{equation}
Indeed, pivoting the five independent singleton coordinates leaves a zero
root row and column and the all-ones ten-pair block. For the last identity,
the five Cayley translations commute and square to the identity, so
$(I+\sum_iT_i)^2=I+5I=0$. The binary code and its optimality among all
storage codes on the base graph are treated in
\cite[Corollary H4]{BrendeCurved}, with Lean proofs there; that paper also
attributes the original binary construction. The cover results below
bound only the specified adjacency code and its subsets.

\subsection{Exact loss under a binary lift}
For a finite connected simple graph $F$, put $C_F=\ker(I+A_F)$ over
$\F_2$. A signing $s:E(F)\to\F_2$ defines the lift with edges
$(u,b)$--$(v,b+s_{uv})$. Let $S_s$ be its symmetric zero-diagonal
signed adjacency matrix and let
\[
 B_s(x,z)=x^TS_sz\quad(x,z\in C_F),\qquad q(s)=\rank B_s.
\]
A signing is a \emph{cut} if $s_{uv}=a_u+a_v$ for vertex values $a_u$.
Call $F$ \emph{signing-rigid} if $B_s=0$ implies that $s$ is a cut.
The reverse implication always holds.

\begin{lemma}[Lift loss and rigidity propagation]\label{lem:binary-loss}
\notformalized
If $k=\dim C_F$, the binary lift $F_s$ satisfies
\[
 \dim C_{F_s}=2k-q(s).
\]
The integer $q(s)$ is even. The lift is connected exactly when $s$ is
not a cut. If $F$ is signing-rigid and $q(s)=2$, then $F_s$ is connected
and signing-rigid.
\end{lemma}
\begin{informalproof}
With $M_F=I+A_F$, the lifted check is
$\left(\begin{smallmatrix}M_F+S_s&S_s\\S_s&M_F+S_s\end{smallmatrix}\right)$.
Writing a word as $(y+z,y)$ gives $M_Fz=0$ and $M_Fy=S_sz$.
Since $\im M_F=C_F^\perp$, the feasible $z$ form $\operatorname{rad}B_s$;
each has a $k$-dimensional fibre of choices of $y$. This proves the
identity. The form is alternating. A nonzero alternating form splits
off a nondegenerate two-dimensional summand; induction proves even rank.

For a cut, $S_s=D_aM_F+M_FD_a$, which vanishes as a form on $C_F$.
Path integration gives the connectivity assertion: if all closed walks
have sign sum zero, their path sums define a vertex potential; otherwise
a closed walk joins the two sheets, and connectedness of $F$ joins all
vertices of the lift.

Suppose now that $F$ is rigid, $q(s)=2$, and a child signing $t$ has
zero form. Pairing diagonal child words $(x,x),(y,y)$ shows that the
sum of its two signs above each base edge has zero base form, hence is
a cut $da$. Switch the child vertices by $a_v$ on sheet zero and by
zero on sheet one. The two signs above each edge now agree, even on
crossed edges, so the switched signing is the pullback of a base
signing $r$. Switching preserves zero form.

Pair a diagonal word with a child word of sheet sum
$z\in\operatorname{rad}B_s$. The resulting value is $B_r(x,z)$.
Every such $z$ occurs, by the lift equations. Thus $B_r$ descends to
the two-dimensional space $C_F/\operatorname{rad}B_s$. Over $\F_2$ its
alternating form is either zero or the form induced by $B_s$.
Rigidity makes $r$ a cut or $s$ plus a cut. Both pull back to cuts:
the pullback of $s$ is the cut given by the sheet-bit potential.
Undoing the first switch proves child rigidity. Since $q(s)>0$, $s$
is not a cut, so the child is connected.
\end{informalproof}

\begin{theorem}[C: sharp binary towers]\label{thm:clebsch-towers}
\notformalized
Start from Clebsch and take any sequence of connected binary lifts.
At a level with $N$ vertices, the binary adjacency code has dimension
at most $N/2+2$. More precisely, either its dimension is $N/2+2$ and
the graph is signing-rigid, or its dimension is at most $N/2$ and stays
at most half at every descendant. Equality $N/2+2$ is attained at
every level by repeatedly signing just one edge. The upper bound also
holds for every connected regular cover of Clebsch with finite
$2$-group deck group, including nonabelian groups.
\end{theorem}
\begin{informalproof}
Clebsch is signing-rigid by \cref{cor:clebsch-rigid} below; that proof
uses only block elimination, not the tower theorem. Write
$e=\dim C_F-|V(F)|/2$. The lift identity gives $e'=2e-q(s)$.
Initially $e=2$. At a rigid parent, a connected lift has positive even
$q(s)$. If $q(s)=2$, the new excess is two and rigidity propagates.
If $q(s)\ge4$, the new excess is nonpositive. Once nonpositive, it
remains so without any rigidity assumption.

Endpoint evaluation of $\ker M$ onto $\F_2^2$ is surjective on every
base edge $uv$. Column $u$ of $M$ has values $(1,1)$ there; column $w$,
for a different neighbor $w$ of $u$, has values $(1,0)$ by triangle
freeness. Both are in the kernel by \eqref{eq:clebsch-base}. Diagonal
kernel words preserve this surjectivity in every binary lift. A
single-edge signing therefore has $q(s)=2$, proving sharpness and
connectedness at every step.

A finite $2$-group has a subgroup chain with successive indices two:
use a central subgroup of order two and induct on the quotient.
The intermediate quotient covers of a connected regular cover remain
connected, and consecutive maps are binary lifts. Normality in the
whole deck group is unnecessary. Apply the tower conclusion.
\end{informalproof}

\subsection{Reciprocal checks and odd characters}
Fix a characteristic-two field $K$ and a positive block width $r$.
A \emph{reciprocal check} $H$ has diagonal $I_r$, zero nonedges of
Clebsch, and edge blocks $G_{uv}\in\mathrm{GL}_r(K)$ with
$G_{vu}=G_{uv}^{-1}$. A vertex gauge replaces each block by
$D_u^{-1}G_{uv}D_v$. The gains are \emph{flat} if every based
closed-walk product is the identity, equivalently if a vertex gauge
makes every edge block $I_r$.

\begin{theorem}[R: reciprocal rank rigidity]\label{thm:clebsch-rank}
\notformalized
Every such check has rank at least $6r$. Equality holds exactly for
flat gains, and rank $6r+1$ is impossible. In particular a scalar
reciprocal check has rank at most seven exactly when it is diagonally
similar to $M$; its rank is then six, and every other one has rank
at least eight.
\end{theorem}
\begin{informalproof}
The full block elimination, including the mixed root row and column,
is given in \cref{sec:clebsch-block-proof}. It proves the lower bound,
the equality criterion, and the forbidden next rank over every
characteristic-two field. The scalar assertion is the case $r=1$.
\end{informalproof}

\begin{corollary}[The base signing rigidity]\label{cor:clebsch-rigid}
\notformalized
The Clebsch graph is signing-rigid.
\end{corollary}
\begin{informalproof}
A binary lift is a reciprocal width-two check, with edge blocks
identity or swap. If $q(s)=0$, \cref{lem:binary-loss} and
\eqref{eq:clebsch-base} give nullity twenty and rank twelve.
By \cref{thm:clebsch-rank} its gains are flat. Every closed walk
therefore has sign sum zero, so path integration makes $s$ a cut.
This argument supplies the base rigidity algebraically, without a
finite rank certificate.
\end{informalproof}

\begin{theorem}[A: direct-product deck groups]\label{thm:clebsch-product}
\notformalized
Let a connected regular Clebsch cover have deck group $P\times D$,
where $P$ is a finite $2$-group and $D$ is a finite odd abelian group.
On its $N=16|P||D|$ vertices the binary adjacency code has dimension
at most $N/2+2$. In particular every finite abelian deck group is
covered. The bound is sharp at every cyclic covering degree.
\end{theorem}
\begin{informalproof}
Choose voltages $(p_{uv},d_{uv})$, with inverse reverse voltages, and
extend scalars to an algebraic closure of $\F_2$. This preserves rank.
The regular representation of the odd abelian direct factor splits
into its $|D|$ characters: commuting translations are diagonalizable,
since their orders are odd, and character orthogonality gives each
character once. The check therefore splits into weighted $P$-cover
checks, with scalar edge weights $\chi(d_{uv})$.

Connectedness says that based closed-walk voltages generate $P\times D$.
Every nontrivial character detects one of their $D$ projections. Its
weighted base check is reciprocal and nonflat, so its nullity is at
most eight by \cref{thm:clebsch-rank}. For any weighted check $B$,
a binary lift with pulled-back weights has equations $Bz=0$, $By=Sz$.
The feasible $z$ form a subspace of $\ker B$, and each fibre has
dimension $\dim\ker B$. Thus nullity at most doubles. A subgroup
chain of $P$ bounds each nontrivial character sector by $8|P|$.
For the trivial character the projected $P$ cover is connected and
unweighted, so \cref{thm:clebsch-towers} instead gives $8|P|+2$.
Summing the sectors gives $8|P||D|+2$. The direct-product assumption
is used in this decomposition; a semidirect product is not included.

For sharpness, put an $m$-cycle permutation $T$ on one base edge $uv$,
its inverse in reverse, and identity elsewhere. Deleting this edge
leaves the base connected, so the cover is connected and regular
cyclic. Put $M_m=M\otimes I_m$, let $U$ embed the endpoint blocks,
and let $V=U^T$. The endpoint surjectivity used in
\cref{thm:clebsch-towers} implies that $V$ is onto on $\ker M_m$.
Symmetry then gives $\im U\cap\im M_m=0$: if $Ua\in\im M_m$, its
pairing with every kernel vector is zero, so surjectivity forces $a=0$.
The changed check is $M_m+U\Delta V$, where
\[
 \Delta=\begin{pmatrix}0&T-I\\T^{-1}-I&0\end{pmatrix}.
\]
Its kernel is exactly $\ker M_m\cap\ker(\Delta V)$, hence has
dimension $10m-\rank\Delta=10m-2\rank(T-I)=8m+2$.
More generally the last term is $8m+2c(T)$, with $c(T)$ the number
of permutation cycles. No diagonalizability of $T$ is needed.
\end{informalproof}

The two bounds give rate at most $1/2+2/N$ for these prescribed checks.
They rule out preserving Clebsch's $5/8$ rate through these cover
operations, even when no separating chain has been specified. They do
not bound arbitrary fitting matrices or nonlinear codes on the lifted
graphs. The tower dichotomy concerns binary towers; it is not asserted
for all direct-product covers. Further small-representation consequences
and a rigid three-sheet seed appear in \cref{sec:clebsch-small}.
No paired-port separation or native implementation follows from any
of these constructions.
